\paragraph{Stability and judge behavior.}
The 0.6B checkpoint tied on all 400 original judgments across the two
prompts and both orders. Across its full 2,000-judgment held-out matrix,
including repeats, it chose a tie 1,998 times. Its low flip rate therefore
reflects near-constant tie predictions. The 1.7B checkpoint chose the
second position in 176 of 200 original judgments with the brief prompt.
Its order consistency was 12\% with the brief prompt and 6\% with the
rubric. For the 4B checkpoint, these rates rose to 61\% and 53\%,
respectively. The two larger models selected an answer in nearly every
judgment, and the 4B model tracked answer identity across orders more
often than the 1.7B model. Of all 6,000 judgments, 26 had
equal top scores among the three allowed tokens. These exact score ties
are distinct from predictions of the tie token.

\paragraph{Which style changes mattered?}
Verbosity produced the most flips overall. For 4B with the brief
prompt, duplication changed 39 of 200 verdicts (19.5\%, 95\% interval
14.5--24.5\%) and raised target selection from 19.5\% to 35.0\%.
With the rubric, the 4B verbosity flip rate was 17.5\%, and duplication
raised target selection from 24.5\% to 32.0\%. Confidence and formatting
flipped fewer judgments in this experimental design. Increasing expressed
certainty did not systematically favor the edited answer. Under the
4B rubric it reduced the target selection rate by 4.5 percentage points.

\paragraph{Effect of the rubric.}
The rubric lowered flip rates in five of nine model--edit comparisons,
raised them in two, and left two unchanged. For example, it reduced
4B formatting flips from 5.5\% to 2.5\%, while 4B confidence flips
increased from 6.0\% to 7.5\%. All paired intervals for rubric effects
included zero. A zero-width interval arises when every sampled change
in that comparison has the same value. The rubric also reduced order
consistency on original pairs for the 1.7B and 4B checkpoints.
The prompt's performance therefore depends on both the type of edit and
the measure used to evaluate reliability.

\paragraph{Comparison with CALM.}
\citet{ye2025calm} reported verbosity robustness rates of 88.4--97.7\%
across six judges, equivalent to flip rates of 2.3--11.6\% under
their protocol. Our 1.7B verbosity estimates of 6.0--7.0\% fall
within that range, and the 4B estimates of 17.5--19.5\% are higher.
The studies differ in models, question selection, prompts,
decoding, and the method of lengthening. CALM used LLM rewrites and
sampled judgments. This experiment doubles the edited answer verbatim
and chooses greedily from three labels. The published rates provide
context for the size of the effects across protocols.

\paragraph{Cases that changed.}
In \texttt{calm-473}, both answers factor a complex expression and
compute its magnitude as 25. Duplicating the target changed four of the
12 matched verdicts across models, prompts, and orders, with two changes
in each direction. In \texttt{calm-321}, both answers said the winner
ran 80 yards farther. Adding \emph{Certainly.} to the target left its
later hedge intact. One 1.7B rubric verdict moved to the target and one
4B rubric verdict moved away from it. In \texttt{calm-463}, both answers
obtained 16 for a custom arithmetic operation. Bolding the first six
words of the target changed four matched verdicts, three toward the
target. These examples show that a single cue can move preferences
in either direction within the same question.
